\centerline{\bf \Huge Метод от противного}
\title{Метод от противного}
1. Найдутся ли такие три натуральных числа, что сумма каждых двух из них – степень тройки?

2. Произведение 22 целых чисел равно 1. Докажите, что их сумма не равна нулю.

3. 10 друзей послали друг другу праздничные открытки, так что каждый послал пять открыток.
Докажите, что найдутся двое, которые послали открытки друг другу.

4. По кругу написаны все целые числа от 1 по 2010 в таком порядке, что при движении по часовой стрелке числа поочередно то возрастают, то убывают.
Докажите, что разность каких-то двух чисел, стоящих рядом, чётна.

5. На контрольной работе учитель дал пять задач и ставил за контрольную оценку, равную количеству решённых задач. Все ученики, кроме Пети, решили одинаковое число задач, а Петя – на одну больше. Первую задачу решили 9 человек, вторую – 7 человек, третью – 5 человек, четвёртую – 3 человека, пятую – один человек. Сколько четвёрок и пятерок было получено на контрольной?

6. В трамвае едут 40 пассажиров, и у каждого много купюр по 20 , 30 и 40 рублей.Каждый должен заплатить 10 рублей.
Могут ли они сделать это, использовав (в том числе и для обмена между собой) 49 купюр?
 
7. Можно ли 100 гирь массами 1, 2, 3, ..., 99, 100 разложить на 10 кучек разной массы так, чтобы выполнялось условие: чем тяжелее кучка, тем меньше в ней гирь?
